Системы управления городским транспортом (UTMS) опираются на данные сенсоров и событий, которые часто оказываются неполными, дублированными или несогласованными по схеме ещё до запуска прогнозирования или обнаружения аномалий. В настоящей работе разрабатывается, реализуется и оценивается \textbf{Flowmatic}~--- интеллектуальная платформа подготовки данных и вывода, которая автоматизирует оценку качества, направляет подготовленные потоки к нейросетевым контрольным точкам соответствующей модальности и публикует воспроизводимые артефакты обучения.

Flowmatic реализует шестимерный индекс качества данных (DQI), автоматизированную очистку и экспорт в контейнеризованном веб-стеке. При пакетной проверке платформы загрузка 30\,000 строк данных Астаны завершается за 406\,мс; при контролируемых искажениях до 20\,\% составной DQI восстанавливается на величину до 0{,}045 при удалении недопустимых строк. Рабочая среда «умного города» объединяет смоделированные источники дорожного движения и погоды, маршрутизацию Manual/Auto по документированному реестру и экспорт в озеро данных типа medallion.

Семь основных нейросетевых контрольных точек обучаются по единому протоколу временного hold-out (разбиение 70/15/15, три seed). Классификатор степени тяжести достигает macro-F1 $= 0{,}911 \pm 0{,}017$, превосходя базовые модели sklearn на той же удерживаемой тестовой выборке; метки представляют собой синтетические атрибуты, заданные правилами. TranAD достигает AUROC $= 0{,}84$ на внедрённых оконных аномалиях. Абляция подготовки снижает RMSE плотности PatchTST с $0{,}90$ до $0{,}58$ после 20\,\% искажения при восстановлении повреждённых ячеек. Правило-ориентированная Auto-маршрутизация на 140 событиях симулятора обеспечивает 100\,\% соответствия сценарной политике при документированных оракулах, а не независимую экспертную точность маршрутизации.

Артефакты опубликованы на платформе Hugging Face в семействе репозиториев \flowmaticrepo{}. Пакетная проверка API выполняется только на экспорте Астаны; развёртывание в реальных ведомствах и внешние метки маршрутизации остаются задачами будущей работы.

\abstractkeywords{Ключевые слова:}{подготовка данных; интеллектуальные транспортные системы; умные города; Flowmatic; маршрутизация моделей; временные ряды; воспроизводимость}
